\chapter{Pendahuluan}

% ================================================================
% PANDUAN PENULISAN BAB I — PENDAHULUAN
% Bab ini menjelaskan alasan dan tujuan penelitian Anda. Tuliskan
% secara naratif, padat, dan didukung sitasi. Perhatikan urutan
% bagian berikut.
% ================================================================

\section{Latar Belakang}

% Tuliskan 4--7 paragraf dengan alur "funnel" (dari umum ke khusus):
%   1) Buka dengan masalah besar/topik yang relevan dengan penelitian.
%   2) Persempit ke masalah spesifik yang akan diteliti.
%   3) Tegaskan urgensi masalah: data prevalensi/insidensi dan dampaknya.
%   4) Tunjukkan kesenjangan (research gap) dari penelitian sebelumnya.
%   5) Akhiri dengan pernyataan singkat maksud penelitian (research aim).
% Setiap pernyataan faktual harus disertai sitasi, contoh:
% \supercite{contohReferensi1,contohReferensi2}
%
% Contoh pola kalimat:
%   "X merupakan masalah yang ...\supercite{ref}. Prevalensinya dilaporkan
%    mencapai ...\supercite{ref}. Namun, data populasi lokal masih terbatas.
%    Oleh karena itu, penelitian ini bertujuan untuk ..."

\section{Rumusan Masalah}

% Uraikan masalah penelitian dalam bentuk pertanyaan penelitian.
% Rumusan masalah umum ditulis dalam satu kalimat tanya; rumusan
% masalah khusus diturunkan menjadi beberapa butir.
%
% Contoh:
% \begin{enumerate}
%     \item Bagaimana karakteristik demografis (usia dan jenis kelamin) subjek penelitian?
%     \item Bagaimana gambaran klinis dan tata laksana subjek penelitian?
% \end{enumerate}

\section{Tujuan Penelitian}

\subsection{Tujuan Umum}

% Nyatakan tujuan penelitian secara luas, selaras dengan judul penelitian.
% Contoh: "Mengetahui karakteristik klinis pasien ..."

\subsection{Tujuan Khusus}

% Rincikan tujuan umum menjadi butir-butir tujuan yang terukur, dengan
% kata kerja operasional (mengetahui distribusi, mengidentifikasi,
% mengukur, menilai, membandingkan, dll).
% \begin{enumerate}
%     \item Tujuan khusus pertama.
%     \item Tujuan khusus kedua.
% \end{enumerate}

\section{Manfaat Penelitian}

% Jelaskan kegunaan hasil penelitian, misalnya:
% - Manfaat teoretis: pengembangan ilmu/konsep di bidang terkait.
% - Manfaat praktis: masukan bagi klinisi, pelayanan kesehatan,
%   dan perencanaan tata laksana atau kebijakan.
% - Manfaat bagi penelitian selanjutnya: dasar/bahan acuan.
